% Venue-agnostic manuscript body: Introduction through Acknowledgments.
% Requires \tool, \previewfig, \adjustwidth and \nameref from the preamble.
\section*{Introduction}

High-throughput technologies for profiling gene expression have revolutionized the field of molecular biology, enabling researchers to capture the intricate patterns of gene regulation and unravel the underlying genetic networks. However, the sheer volume and complexity of gene expression data pose significant challenges in extracting meaningful insights \cite{lorenaComplexityGeneExpression2008}. One critical step in this process is the clustering of genes based on their expression profiles, which aims to group together genes exhibiting similar patterns of expression across various experimental conditions or time points \cite{Saelens2018}. Clustering is a foundational step in computational biology, from grouping co-expressed genes into putative regulatory modules~\cite{Eisen1998,Wolfe2005} to shedding light on regulatory pathways mechanisms \cite{vengatharajulooGeneCoExpressionNetwork2023}, disease-associated genes \cite{vandamGeneCoexpressionAnalysis2017}, and potential therapeutic targets \cite{hasankhaniDifferentialCoExpressionNetwork2021}. By partitioning a set of objects into groups based on similarity, cluster analysis identifies intrinsic structures in data to gain insight \cite{hruschka2004, pirim2012}.  

The dominant practice has two stages. First, a fixed algorithm such as k-means~\cite{MacQueen1967} or hierarchical clustering~\cite{Eisen1998} produces a partition. Second, a cluster-validity index like the silhouette~\cite{Rousseeuw1987}, 
the Davies--Bouldin index~\cite{DaviesBouldin1979}, the Calinski--Harabasz
criterion~\cite{CalinskiHarabasz1974}, or an information criterion such as
BIC~\cite{Schwarz1978} is used to \emph{evaluate} that partition, often
to choose the number of clusters. The objective the algorithm optimizes (for
k-means, within-cluster sum of squares) and the objective by which the result is
ultimately judged are therefore distinct. When the two disagree, the reported
partition is optimal for neither the user's stated criterion nor, necessarily,
the underlying biology.

This gap is particularly consequential in gene co-expression analysis. There,
the notion of a good module is intrinsically correlational: genes in a module
should be strongly co-expressed across conditions, regardless of their absolute
expression levels or Euclidean geometry. Tools such as WGCNA~\cite{Langfelder2008},
Clust~\cite{AbuJamous2018}, and UNCLES~\cite{AbuJamous2015} encode this intuition
through correlation networks or consensus procedures, and comparative studies
show that module-detection performance depends strongly on how the objective is
framed~\cite{Saelens2018}. Yet the correlation-based quality of the final
modules is still, in most pipelines, assessed only after the clustering
algorithm has finished optimizing a different, geometric objective.

We argue for inverting the workflow by making the validity metric itself the
optimization target. Concretely, the user specifies the scalar criterion that
defines a good partition for their problem, and a search procedure evolves a
partition to optimize that criterion directly. This is attractive whenever the
criterion of interest is (i) not what standard algorithms optimize, (ii)
non-differentiable or combinatorial, or (iii) most naturally expressed as a
trade-off between competing concerns---fit versus parsimony, compactness versus
separation. Direct optimization of clustering criteria has a long history in
evolutionary computation~\cite{Maulik2000,Handl2007}, but it is rarely packaged
as a practical, dependency-light tool aimed at the co-expression use case and
built around a correlation objective.

While traditional numerical optimizers---such as Newton-Raphson, Powell's method, or the Simplex 
algorithm---are highly effective for continuous parameter estimation, they are fundamentally 
ill-suited for this clustering framework. These methods require a continuous search space and 
differentiable objective functions to calculate valid step directions. In contrast, identifying 
optimal partitions is an inherently discrete, combinatorial problem. Our objective functions are 
highly discontinuous with respect to categorical cluster labels, resulting in a rugged fitness 
landscape riddled with local optima where gradient-based methods become trapped. Furthermore, 
while stochastic Bayesian approaches (such as Reversible Jump MCMC) could theoretically navigate
 this discrete landscape by jumping between partition states, the computational time cost of 
 sufficiently sampling such a vast combinatorial space is prohibitive for large datasets. 
 The Genetic Algorithm provides the necessary robustness to directly evolve discrete partition 
 vectors, efficiently navigating the non-differentiable landscape via crossover and mutation without 
 the extreme computational overhead of exhaustive Bayesian sampling or the restrictive assumptions 
 of gradient-based methods.

Here we present \tool{} (Genetic Clustering by Metric), an open-source Python package that clusters data by directly optimizing a chosen cluster-validity metric with a memetic genetic algorithm (GA), which is an evolutionary algorithm hybridized with a problem-specific local search that refines each candidate solution within its own ``lifetime''. This strategy was introduced by Moscato and later surveyed by Neri and Cotta~\cite{Moscato1989,NeriCotta2012}. Genetic algorithms (GAs) could be a promising solution for clustering tasks \cite{katoch2021}, particularly in the context of gene expression data \cite{pirim2012}. Other attempts have been made to use evolutionary algorithms in this or similar contexts, with overall promising results \cite{hageman2008, hruschka2006, korkmaz2006, weiEntropyClusteringAnalysis2008}. GAs are optimization algorithms that mimic the principles of natural selection and evolution, making them well-suited for complex, multi-objective problems where traditional methods may struggle to find optimal solutions \cite{wirsanskyHGA2020}. 

\tool{} makes four contributions. First, we place co-expression module detection
on an explicit likelihood framework: the absolute-correlation criterion used
throughout the field is the profile log-likelihood of a block-diagonal one-factor
Gaussian with equal-magnitude loadings, and it is the constrained member of a
family whose general form permits $q$ factors with free loadings. Making that
nesting visible turns an implicit modelling choice---that a module's members
respond to their regulator equally strongly, and to one regulator only---into a
testable assumption with a strictly more general alternative. Second, we hybridize
the GA with an efficient memetic local search and benchmark it against the
community-detection and network methods that dominate practice---Leiden
\cite{Traag2019}, Louvain~\cite{Blondel2008}, spectral
clustering~\cite{NgJordanWeiss2002}, Markov clustering~\cite{VanDongen2008} and
WGCNA~\cite{Langfelder2008}---each given the true number of modules or tuned
against ground truth, so that no margin can be attributed to unfavourable settings
for a competitor. Third, we test on real data with real ground truth, using an established module-detection benchmark~\cite{Saelens2018}, which pairs expression compendia
from three organisms with curated module definitions derived from regulatory
networks. Fourth, we characterize when our own
objective is wrong: across data-generating processes that violate the block model
we identify two structural regimes in which maximizing the constrained likelihood
moves \emph{away} from the true modules, demonstrate that this is model
misspecification rather than optimizer failure, and show that the free-loading
model removes the failure on those processes. The implementation is deliberately minimal---depending only on NumPy~\cite{Harris2020}, SciPy~\cite{Virtanen2020}, pandas~\cite{McKinney2010} and Matplotlib~\cite{Hunter2007}, with internal
k-means++~\cite{ArthurVassilvitskii2007} and correlation-distance hierarchical
seeders, behind one swappable-metric interface that also exposes standard
geometric indices and a multi-objective (NSGA-II~\cite{Deb2002}) mode.

\section*{Materials and methods}

\subsection*{Overview and problem statement}
\tool{} clusters $n$ \emph{elements} (the rows of a data matrix
$\mathbf{X}\in\mathbb{R}^{n\times d}$; in co-expression analysis, $n$ genes
measured across $d$ samples) by searching the space of partitions for one that
optimizes a user-chosen validity metric. A partition is represented as an
integer label vector $\boldsymbol{\ell}\in\{0,1,\dots,K\}^{n}$, where
$\ell_i = c$ assigns element $i$ to cluster $c$ and the optional label $0$ marks
an element as unassigned. The user fixes an upper bound $g_{\max}$ on the number
of clusters. Each candidate partition $\boldsymbol{\ell}$ is scored by one or
more metric functions $m(\boldsymbol{\ell})\in\mathbb{R}$; the GA evolves a
population of partitions to optimize these scores in their natural direction
(maximize or minimize).

\subsection*{Performance-target metrics}
\tool{} treats every metric as an interchangeable ``performance target'' or fitness measure \cite{wirsanskyHGA2020}. Table~\ref{tab:metrics} lists the built-in targets, the input format each one is designed to work with, and its optimization direction. They fall into two families: correlation-based objectives designed for modular data, and standard geometric/probabilistic indices.

\paragraph*{A generative model for correlation modules.}
The correlation objective is the maximum log-likelihood of an explicit
generative model, which we state here and derive in full in
\nameref{AppendixA}. Let the data matrix $\mathbf{X}\in\mathbb{R}^{n\times d}$
hold $n$ genes measured over $d$ samples, row-standardized. We model each
\emph{sample} (column) as an independent draw from a multivariate Gaussian over
genes, $N(\mathbf{0},\Sigma)$, whose correlation matrix $\Sigma$ is
\emph{block-diagonal} by the partition $\boldsymbol{\ell}$: genes in the same
module are correlated, genes in different modules independent. Within a module
$s$ we adopt the canonical single-regulator (one-factor) structure with
equal-magnitude loadings of arbitrary sign, $x_g = \sigma_g f_s + \varepsilon_g$
with $\sigma_g\in\{\pm\lambda_s\}$, so that any two genes in the module share
$\lvert\mathrm{corr}\rvert = \rho_s$. This is exactly the core assumption of a
co-expression module (one regulator, up- or down-regulating its targets).

\paragraph*{The likelihood reduces to the project's objective.}
Let $\mathbf{R}$ be the Pearson correlation matrix between genes and, for module
$s$ of size $n_s$, let
\begin{equation}
\label{eq:cs}
C_s = \sum_{i\in s}\sum_{j\in s} \lvert R_{ij}\rvert \in [\,n_s,\,n_s^{2}\,]
\end{equation}
(the unit diagonal is included). As shown in \nameref{AppendixA}, the
maximum-likelihood common coherence of a module is its mean absolute
off-diagonal correlation, $\hat\rho_s=(C_s-n_s)/(n_s^2-n_s)$, and at that
estimate the trace term of the Gaussian log-likelihood collapses to the module
dimension. The profile log-likelihood of a partition is therefore, up to a
partition-independent constant,
\begin{equation}
\label{eq:loglik}
\ell(\boldsymbol{\ell}) = d\,\mathcal{L}(\boldsymbol{\ell}),
\qquad
\mathcal{L}(\boldsymbol{\ell}) = \frac{1}{2}\sum_{s\,:\,n_s>1}
\left[\, \log\!\frac{n_s}{C_s}
   + (n_s-1)\,\log\!\frac{n_s^{2}-n_s}{\,n_s^{2}-C_s\,} \,\right],
\end{equation}
where $\mathcal{L}$ is precisely $-\tfrac12\sum_s\log\det\Sigma_{\mathrm{CS}}(\hat\rho_s)$,
the equicorrelation log-determinant. Maximizing $\mathcal{L}$ thus maximizes the
model likelihood. Each summand is monotone in module coherence: as
$C_s\to n_s^2$ the cluster is perfectly correlated (best fit) and as $C_s\to n_s$
it is incoherent (contributing zero); singletons contribute nothing. A perfectly
correlated cluster is awarded a large finite reward rather than $+\infty$ for
numerical stability.

\paragraph*{Model selection: information criteria.}
The raw likelihood favours more modules, since splitting a coherent module into
coherent sub-modules does not necessarily lower $\mathcal{L}$. Because each module carries a single coherence parameter $\rho_s$, the model with $K$ modules has $K$ free
parameters, and the standard information criteria over the $d$ samples are
\begin{equation}
\label{eq:loglikic}
\mathrm{BIC}(\boldsymbol{\ell}) = -2 d\,\mathcal{L}(\boldsymbol{\ell}) + K\log d,
\qquad
\mathrm{AIC}(\boldsymbol{\ell}) = -2 d\,\mathcal{L}(\boldsymbol{\ell}) + 2K,
\end{equation}
to be minimized (targets \texttt{loglik\_bic}, \texttt{loglik\_aic}). The
likelihood scales with the sample size $d$ while the penalty does not, so more
samples correctly license more modules. This is distinct from the geometric
\texttt{bic}/\texttt{aic} targets, whose Euclidean-Gaussian assumption conflicts
with the absolute-correlation objective when modules contain anti-correlated
members, which is precisely the regime in which the correlation criteria are the right
choice.

\paragraph*{Geometric and probabilistic indices.}
For data whose clusters are defined geometrically, \tool{} provides four
standard indices computed on $\mathbf{X}$: the mean silhouette
coefficient~\cite{Rousseeuw1987} (maximize; range $[-1,1]$), the Davies--Bouldin
index~\cite{DaviesBouldin1979} (minimize), the Calinski--Harabasz
variance-ratio criterion~\cite{CalinskiHarabasz1974} (maximize), and the
Bayesian and Akaike information criteria~\cite{Schwarz1978,Akaike1974} under a
spherical-Gaussian-per-cluster model with pooled variance (both minimize). The
BIC/AIC log-likelihood follows the x-means formulation~\cite{Pelyhe2000}; for a
partition with $k$ clusters in $d$ dimensions the model has
$p = kd + (k-1) + 1$ free parameters (centroids, mixture weights, and one pooled
variance), and BIC $=-2\log\hat{L} + p\log n$, AIC $=-2\log\hat{L} + 2p$. Unlike
the correlation objective, BIC and AIC reward parsimony explicitly and so can be
used for model selection on geometrically compact data.

\begin{table}[!ht]
\centering
\caption{\textbf{The built-in performance-target metrics.}}
\label{tab:metrics}
\begin{tabular}{@{}llll@{}}
\hline
Target & Direction & Input & Role \\
\hline
\texttt{loglik}            & maximize & correlation matrix & correlation coherence (modules) \\
\texttt{loglik\_bic}       & minimize & correlation matrix & coherence with cluster-count penalty \\
\texttt{loglik\_aic}       & minimize & correlation matrix & coherence with lighter penalty \\
\texttt{loglik\_factor@}$q$ & maximize & correlation matrix & rank-$q$ free-loading factor fit \\
\texttt{loglik\_factor\_bic@}$q$ & minimize & correlation matrix & factor fit with parameter penalty \\
\texttt{loglik\_factor\_aic@}$q$ & minimize & correlation matrix & factor fit, lighter penalty \\
\texttt{silhouette}        & maximize & data matrix        & geometric separation \\
\texttt{davies\_bouldin}   & minimize & data matrix        & compactness/separation ratio \\
\texttt{calinski\_harabasz}& maximize & data matrix        & variance-ratio criterion \\
\texttt{bic}               & minimize & data matrix        & Gaussian-mixture parsimony \\
\texttt{aic}               & minimize & data matrix        & Gaussian-mixture parsimony \\
\hline
\end{tabular}
\begin{flushleft}
The correlation-based targets consume the element-by-element Pearson correlation
matrix; the geometric/probabilistic indices are computed directly on the
(standardized) data matrix. Any subset can be combined under multi-objective
optimization. The penalized correlation criteria \texttt{loglik\_bic} and
\texttt{loglik\_aic} add a cluster-count penalty to \texttt{loglik} so that a
single objective both fits and selects the number of modules (see below). The factor targets take their rank $q$ as a suffix (e.g.\ \texttt{loglik\_factor@3}); their penalized forms select the rank but not the number of modules (Results).
\end{flushleft}
\end{table}

\subsection*{Data preparation}
Input is a numeric matrix (CSV/TSV or array) with rows as elements and columns
as features; header rows and index columns are auto-detected. By default the
data are standardized to zero mean and unit variance. For the correlation
objective, standardization is applied per row (per element/gene), which is the
regime under which the Pearson correlation between rows is meaningful; for the
geometric indices, per-feature standardization is the appropriate choice. The
element-by-element correlation matrix $\mathbf{R}$ is computed once and reused
across all evaluations of $\mathcal{L}$.

\subsection*{Dimensionality-reduction visualization}
Visual inspection of clustering results is essential for communicating and
sanity-checking a partition.  \tool{} provides built-in tools for projecting the
clustered data into two (or three) dimensions, with the cluster labels shown as
colours.  Two embedding methods are always available (no extra dependencies):
\emph{principal component analysis} (PCA), a linear projection preserving maximal
variance, and \emph{correlation multidimensional scaling} (correlation MDS), a
classical MDS embedding~\cite{Torgerson1952} on the correlation distance matrix
$1 - \lvert\mathbf{R}\rvert$.  The second method is particularly informative for
the correlation objective because it projects the elements into the same distance
space the log-likelihood operates in: if modules are well-separated in
correlation, they appear well-separated in the MDS embedding.  For users with
additional packages installed, t-SNE~\cite{vanderMaaten2008} and
UMAP~\cite{McInnes2018} are also supported and may better resolve local
neighbourhood structure in large or high-dimensional datasets.  These tools are
exposed both through the Python API (\texttt{reduce\_dimensions} and
\texttt{plot\_clusters}) and the command line (\texttt{gcmrk visualize}).

\subsection*{Genetic algorithm}
\tool{} evolves a population of partitions with a self-contained GA. The components are:

\begin{enumerate}
\item \textbf{Initialization.} The initial population is seeded with three kinds
of partition across candidate cluster counts $k = 2,\dots,g_{\max}$: k-means++
partitions (a small internal NumPy k-means with multiple restarts);
\emph{correlation-aware} partitions obtained by agglomerative (average- and
complete-linkage) hierarchical clustering of the correlation distance
$1-\lvert\mathbf{R}\rvert$, the classical co-expression clustering strategy; and
random partitions for diversity. User-supplied seeds (e.g.\ the output of WGCNA
or Clust) can also be injected. The correlation-aware seeds are important for the
\texttt{loglik} objective: Euclidean k-means splits strongly anti-correlated
elements that are coherent under the absolute-correlation criterion, whereas the
correlation-distance seeds already respect the structure the GA is asked to
optimize.
\item \textbf{Evaluation.} Each partition is scored on the chosen target(s); a
cache keyed on the label vector avoids recomputing identical partitions.
\item \textbf{Selection and variation.} In \emph{weighted} mode the targets are
combined into a single fitness through direction-aware weights after a bounded
squashing transform that places metrics of different scales on a comparable
footing, and elitist tournament selection is applied. In \emph{nsga2} mode the
targets are treated as competing objectives and NSGA-II~\cite{Deb2002}
fast non-dominated sorting with crowding distance is used. Variation is two-point
crossover on label vectors plus a reassignment mutation that moves a small
fraction of elements to random clusters.
\item \textbf{Repair.} After every operator a partition is relabelled to
consecutive integers and constrained to have between $2$ and $g_{\max}$ clusters,
so that all validity metrics remain well defined (degenerate single-cluster or
empty-cluster partitions are repaired).
\item \textbf{Memetic local search.} For the correlation objective, the seed
partitions and the elite of each generation are refined by a greedy hill-climb
that repeatedly applies the single element reassignment most improving the
objective until none remains (a memetic, Lamarckian GA). The gain of a move
touches only the source and target modules, so it is computed incrementally in
$\mathcal{O}(\text{module size})$ from $C_s$ (Eq~\eqref{eq:cs}); a full sweep is
$\mathcal{O}(n^2)$. This step performs exactly the move that agglomerative clustering---used here only for seeding---cannot, namely reassigning an element after an early merge; the ablation in Results shows it accounts for much of the optimizer's margin at moderate noise, though not at the highest noise level tested. When the objective is a penalized
criterion (Eq~\eqref{eq:loglikic}) the hill-climb optimizes that criterion, so it
also merges spurious modules.
\end{enumerate}

For multi-objective runs \tool{} returns the full Pareto front and, as a single
representative partition, the knee point that maximizes the sum of min--max
normalized objectives across the front. All experiments use crossover probability 0.8, mutation probability 0.2 with per-gene reassignment probability 0.05, tournament size 3, elitism 5 and two k-means restarts per candidate $k$, with a fixed random seed. Population size and generations are set per experiment: a population of 150 evolved for 100 generations for the comparative benchmark, the misspecification sweep, the scaling study and the diagnostic experiments; 100 for 60 generations for the two real-data analyses, where the factor objective is far more costly; and 200 for 120 generations for the single-dataset illustrations on the synthetic designs and Iris. Within any one comparison every \tool{} variant uses the same settings.

\subsection*{Choosing the number of clusters}
Because $\mathcal{L}$ (Eq~\eqref{eq:loglik}) has no parsimony term (penalization against more clusters), \tool{} offers three routes to the right number of clusters. (i) \emph{Constrain} $g_{\max}$ to the expected number of modules. (ii) \emph{Optimize a penalized correlation criterion} (\texttt{loglik\_aic} or \texttt{loglik\_bic},
Eq~\eqref{eq:loglikic}) as a single objective under a loose $g_{\max}$; this
selects $K$ within correlation space and is our recommended default when $K$ is
unknown. Equivalently, one can scan $g_{\max}$ with \texttt{loglik} and pick the
$K$ that minimizes Eq~\eqref{eq:loglikic}. (iii) \emph{Explore an explicit
trade-off} by optimizing \texttt{loglik} jointly with a parsimony target under
NSGA-II and reading the knee of the Pareto front; pairing with the geometric
\texttt{bic} is informative only when the modules are also geometrically compact,
whereas the penalized correlation criteria avoid that assumption.

\subsection*{The factor block model}
\label{sec:factor}
Eq~\eqref{eq:loglik} enters a module only through the sum of absolute
correlations inside it. Every within-module pair is therefore exchangeable: only
the \emph{mean} within-module $\lvert r\rvert$ matters, not how it is distributed
across members. That is precisely the equal-magnitude $\pm$ loading assumption,
and it is restrictive in two ways that matter biologically---members of a real
module respond to their regulator with different strengths, and a module may
answer to more than one regulator. Relaxing both gives a strictly more general
model, of which Eq~\eqref{eq:loglik} is a constrained special case; whether the
extra freedom pays is an empirical question we answer in Results, and the answer
differs between synthetic and real data.

For a module $m$ with $n_m$ genes we take the block of the correlation matrix to
follow a $q$-factor structure with free loadings,
\begin{equation}
\mathbf{R}_m \;=\; \boldsymbol{\Lambda}_m\boldsymbol{\Lambda}_m^{\!\top} + \sigma_m^2\mathbf{I},
\qquad \boldsymbol{\Lambda}_m \in \mathbb{R}^{n_m \times q},
\label{eq:factorblock}
\end{equation}
which is probabilistic PCA~\cite{TippingBishop1999} applied blockwise. Its
maximum-likelihood fit is available in closed form from the eigenvalues of
$\mathbf{R}_m$: with $\ell_{m,1}\ge\dots\ge\ell_{m,n_m}$ the eigenvalues of the
block, the residual variance is the mean of the discarded eigenvalues,
$\hat{\sigma}^2_m = (n_m-q)^{-1}\sum_{i>q}\ell_{m,i}$, and the profile
log-likelihood contribution of the module is
\begin{equation}
\mathcal{L}^{\text{factor}}_m \;=\; -\frac{d}{2}\left[\sum_{i\le q}\log \ell_{m,i}
\;+\;(n_m-q)\log\hat{\sigma}^2_m\right].
\label{eq:factorloglik}
\end{equation}
Only the leading $q$ eigenvalues and the block trace are needed, which keeps the
objective affordable inside the genetic algorithm. Two properties matter. First,
Eq~\eqref{eq:factorblock} is fitted to the \emph{signed} correlation matrix: a
factor with mixed-sign loadings reproduces an anti-correlated module natively, so
the absolute value that Eq~\eqref{eq:loglik} requires is no longer needed.
Second, the number of free parameters,
$p_m = n_m q - q(q-1)/2 + 1$, grows with module size, where the exchangeable
model charges one parameter per module regardless of size. Summing $p_m$ into the
usual AIC/BIC penalties therefore yields criteria that are comparable
\emph{across ranks}, so $q$ need not be fixed by the user: we run the optimizer
once per candidate rank and keep the fit minimizing BIC. We use BIC rather than
AIC throughout, because AIC selects one rank too many (\nameref{AppendixA}).

\emph{Relationship between the two objectives.} Eq~\eqref{eq:factorblock} at
$q=1$ with the loadings constrained to a common magnitude,
$\lvert\lambda_{mi}\rvert=\lambda_m$ for all $i$ in module $m$, reproduces the
exchangeable model of Eq~\eqref{eq:loglik}: under that constraint every
within-module correlation has the same absolute value $\lambda_m^2$, which is why
Eq~\eqref{eq:loglik} can work with $\lvert\mathbf{R}\rvert$ and a single
per-module coherence parameter. The two are therefore not competing objectives
but nested models, and Eq~\eqref{eq:loglik} is the fast, strongly constrained
member of the family: it requires no rank selection and is dramatically cheaper
in practice---on real compendia we measured the factor procedure at 20 to 110
times the cost of the constrained criterion, because real correlation structure
yields more modules and a higher selected rank than synthetic designs do---at the
price of the two assumptions above. We report both throughout, and quantify what
the constraint costs and what relaxing it buys.

Because the factor likelihood does not decompose into the $O(1)$ incremental
updates the exchangeable objective admits, the local search re-solves the
eigendecomposition of affected blocks and screens candidate destinations to the
three most similar modules; on inputs above $1500$ elements the refinement step is
skipped rather than allowed to dominate the run.

\subsection*{Competing methods and the fairness protocol}
\label{sec:baselines}
Because our central question is comparative, every competitor is given the most
favourable setting we can construct. Methods that accept a cluster count---
k-means, spectral clustering~\cite{NgJordanWeiss2002} on the affinity
$\lvert\mathbf{R}\rvert^{\beta}$, and agglomerative clustering on
$1-\lvert\mathbf{R}\rvert$---are handed the true $K$. Leiden~\cite{Traag2019} and
Louvain~\cite{Blondel2008} take a resolution rather than a cluster count, and
Markov clustering~\cite{VanDongen2008} an inflation parameter; for the known-$K$
comparison we bisect these parameters until each method returns exactly $K$
communities, which is the most generous setting available to them. We verified
that the graph methods are insensitive to the soft-thresholding power $\beta$
(bisecting to $K$ renormalizes granularity), so the shared choice $\beta=6$ does
not handicap them.

WGCNA~\cite{Langfelder2008} is reported twice, because it selects its own module
count from a fixed tree-cut height and is correspondingly sensitive to
thresholding. The \emph{default} arm runs the standard pipeline (soft threshold
chosen by \texttt{pickSoftThreshold}, topological overlap, dynamic hybrid tree
cut~\cite{LangfelderHorvath2008DynamicTreeCut}), which is what a practitioner
following the tutorial obtains. The \emph{oracle-tuned} arm reports the best
result over a grid of soft-thresholding powers and minimum module sizes,
\emph{scored against the ground truth}---an advantage no real user has. We
compare against the oracle arm so that no margin we report can be attributed to
unlucky WGCNA hyper-parameters.

\subsection*{Data-generating processes that violate the model}
\label{sec:misspec}
Recovering modules from \tool{}'s own simulator demonstrates that the optimizer
reaches the likelihood optimum; it cannot demonstrate that the model describes
real co-expression data, since the simulator and the objective encode the same
assumptions. We therefore generate data from five further processes, each
breaking a different assumption while retaining planted ground truth.
\emph{Multi-factor} modules are driven by three latent factors with continuous
Gaussian loadings, so within-module correlation is heterogeneous and some
same-module pairs are nearly independent. \emph{Hub-and-spoke} modules carry a
single factor whose loading decays geometrically across members, so members are
not exchangeable. \emph{Overlapping} modules give a fraction of genes a secondary
loading on a second module, so no hard partition is exactly correct.
\emph{Negative-binomial counts} replace the Gaussian with a gamma-Poisson
mixture on the log scale, breaking distributional assumptions and adding the
mean--variance trend of RNA-seq. \emph{Network} data are drawn from a Gaussian
whose covariance arises by diffusion over an LFR benchmark
graph~\cite{LancichinettiFortunato2008} with planted communities and power-law
degrees, so correlation reflects connectivity through a sparse network rather
than a shared factor. Each generator's noise level is calibrated so that the
within- minus between-module absolute-correlation contrast is comparable across
processes; we report that contrast alongside every result, because a generator
whose contrast has collapsed cannot discriminate between methods.

\subsection*{Benchmark datasets and evaluation}
We evaluate on synthetic designs with planted modules, on the misspecified processes described above, on Iris, on real expression compendia with curated modules, and on a clinical RNA-seq dataset.

\emph{Synthetic-Easy} and \emph{Synthetic-Hard} are generated by \tool{}'s
modular simulator: each module is driven by a shared latent factor (with random
$\pm 1$ loadings, so within-module genes are strongly correlated in absolute
value while genes in different modules are not), with per-entry Gaussian noise
added and rows finally standardized. Synthetic-Easy has three equal modules
(\EasyGenes{} genes across \EasySamples{} samples) at low noise, giving strong
separation (mean within-module $\lvert r\rvert = \EasyWithinR{}$ versus
between-module $\EasyBetweenR{}$). Synthetic-Hard has five unequal modules
(15, 20, 25, 30, 10 genes; \HardGenes{} genes across \HardSamples{} samples) at
higher noise, giving weak separation (within $\HardWithinR{}$ versus between
$\HardBetweenR{}$).

\emph{Iris} \cite{Fisher1936} is the classic \IrisSamples{}-sample,
\IrisFeatures{}-feature, 3-class benchmark; it is low-dimensional and geometric
rather than correlational, and we use it to show that the same engine, driven by
geometric targets, faithfully optimizes those indices on data outside the
co-expression regime.

\emph{Curated-module benchmark.} For real data with ground truth we use an established module-detection benchmark~\cite{Saelens2018}, from which we take \SaelensNDatasets{} expression compendia spanning three organisms---\emph{E. coli} (COLOMBOS, PRECISE2), yeast (GPL2529, DREAM5) and human (TCGA, GTEx)---with module definitions derived from regulatory networks: RegulonDB regulons for \emph{E. coli}, MacIsaac and Sisima regulons for yeast, and regulatory circuits for human. We analyse the \SaelensNTop{} most variable genes of each compendium, intersect the known modules with that gene set, and discard known modules left with fewer than five genes or spanning more than a quarter of the analysed genes. The human regulatory circuits are published at nine stringency thresholds; we use the most stringent, because at permissive thresholds a single ``module'' spans a quarter of the transcriptome. Regulons overlap and do not cover every gene, so partitions are scored not by ARI but by the benchmark's own recovery and relevance measures (Results).

\emph{Clinical RNA-seq data.} We also analyse GSE183947~\cite{Zhang2021}, a breast-cancer RNA-seq dataset of \EmpSamples{} samples (30 tumours and 30 patient-matched adjacent-normal tissues) with \EmpFullGenes{} genes quantified as FPKM. After $\log_2$ transformation every method is applied to the \EmpCompNTop{} most variable genes, and the partitions are scored by within-module coherence; by the fraction of modules with Gene Ontology Biological Process over-representation (Enrichr, with the analysed genes as background); by whether module eigengenes separate tumour from normal tissue (Welch $t$-test, Benjamini--Hochberg); and by split-half reproducibility, in which the samples are divided in half within each tissue type, each half is clustered independently and the two partitions are compared by ARI. None of these criteria uses the tissue label to fit anything.

\emph{Comparative benchmark.} To test whether globally optimizing the block
likelihood improves on the alternatives a practitioner would otherwise use, we
run a replicated benchmark on the five-module design over a noise sweep
$\sigma\in\{0.7,1.0,1.5,2.0\}$, with \NSeeds{} random datasets per noise level.
We compare \tool{} (memetic) and \tool{} driven by the factor objective with
BIC-selected rank, against: \tool{} with the local-search step disabled (an
ablation isolating its contribution); agglomerative clustering on
$1-\lvert\mathbf{R}\rvert$ with average and with complete linkage, cut at the
true $k$; k-means; spectral clustering; Leiden; Louvain; Markov clustering; and
WGCNA in both its default and oracle-tuned arms, under the fairness protocol
described above. We also evaluate model selection with $k$ unknown---where every
method runs free, with no ground-truth $k$ and no oracle tuning---and robustness
to unstructured ``noise'' genes that belong to no module. Because every method
sees the same datasets, comparisons are paired, and we test them with the
Wilcoxon signed-rank test across the \NSeeds{} replicates rather than comparing
marginal means.

\emph{Scaling.} We measure wall-clock runtime and recovery on the modular design
at $200$, $500$, $1000$, $2000$ and $\ScaleMaxGenes{}$ genes with ten modules,
three replicates per size, on the same hardware and with identical GA settings.

We use the adjusted Rand index (ARI)~\cite{HubertArabie1985}; for the single-dataset summaries we also report Hungarian-matched accuracy. Results are reported as mean $\pm$ SD over the random datasets. All code to reproduce the experiments, figures, and tables is provided with the package.

\section*{Results}

\subsection*{Recovering correlated modules on synthetic data}
On Synthetic-Easy, \tool{} driven by the correlation log-likelihood (with
$g_{\max}$ set to the true count) recovers the planted three-module structure
perfectly (ARI $=\EasyLoglikARI{}$; Table~\ref{tab:benchmark}), whereas the k-means baseline, given the
same true $k$, fails badly (ARI $=\EasyKmeansARI{}$). The gap is diagnostic:
because module loadings carry random signs, each module contains strongly
anti-correlated genes that are coherent under the absolute-correlation objective
but lie far apart in Euclidean space, so k-means splits them. The GA converges
within the first few dozen generations (Fig~\ref{fig:convergence}). On
Synthetic-Hard, where weak separation and higher noise make the problem
substantially harder, \tool{} still recovers the five modules well
(ARI $=\HardLoglikARI{}$) while k-means remains poor (ARI $=\HardKmeansARI{}$).
The correlation-aware hierarchical seeding (Methods) is important here: it
supplies the GA with starting partitions that already respect the
absolute-correlation structure, which a Euclidean seeder cannot.

% Place tables after the first paragraph in which they are cited.
\begin{table}[!ht]
\centering
\caption{\textbf{Benchmark results: clusters recovered ($k$), adjusted Rand
index (ARI), and Hungarian-matched accuracy against ground truth.}}
\label{tab:benchmark}
\begin{adjustwidth}{-1.25in}{0in}
\centering
\IfFileExists{results/benchmark_table.tex}{\input{results/benchmark_table.tex}}{%
\begin{tabular}{@{}l@{}}\hline
\textit{Run experiments/run\_experiments.py to populate this table.}\\\hline
\end{tabular}}
\end{adjustwidth}
\begin{flushleft}
True $k$ is 3 (Synthetic-Easy, Iris) and 5 (Synthetic-Hard); GSE183947 has no
ground-truth labelling (ARI/accuracy not applicable). ``known $k$'' fixes
$g_{\max}$ to the true value; ``loose $g_{\max}$'' and ``model sel.'' use a loose
bound, the latter with a penalized correlation criterion. k-means is given the
true $k$ and 10 restarts. GA settings as in Methods; fixed random seed.
\end{flushleft}
\end{table}

\begin{figure}[!h]
\caption{\textbf{Convergence of the genetic algorithm on Synthetic-Easy.}
Best and population-mean weighted fitness (correlation log-likelihood) across
generations. The optimizer reaches a high-quality partition within the first few
dozen generations and the elite fitness is monotone non-decreasing by
construction.}
\previewfig{convergence.pdf}
\label{fig:convergence}
\end{figure}

\subsection*{Comparison against contemporary module detectors}
The central methodological question is whether globally optimizing the block
likelihood improves on the methods a practitioner would otherwise reach for. The
replicated noise sweep (Fig~\ref{fig:perfnoise}, Table~\ref{tab:baselines}) shows
that it does at every noise level where the methods can be distinguished, though the size of the margin and which
variant delivers it both depend on the noise.

Driven by the constrained objective of Eq~\eqref{eq:loglik}, \tool{} reaches ARI
$\MemeticHardARI{}$ at $\sigma=1.5$, ahead of every competitor and significantly
so under a paired Wilcoxon signed-rank test across the \NSeeds{} shared datasets
($p<0.001$ against each): $\WgcnaOraHardARI{}$ for oracle-tuned WGCNA,
$\SpectralHardARI{}$ for spectral clustering, $\LeidenHardARI{}$ for Leiden,
$\LouvainHardARI{}$ for Louvain and $\HierAvgHardARI{}$ for average-linkage
hierarchical clustering at the true $K$. At lower noise every method except
k-means and MCL saturates near ARI $1.0$ and the comparison is uninformative.

The one place the constraint costs something here is at the highest noise level:
at $\sigma=2.0$ the constrained objective falls to $\MemeticExtremeARI{}$ and
oracle-tuned WGCNA overtakes it (paired difference $-0.11$, $p=1\times10^{-4}$),
whereas the free-loading model of Eq~\eqref{eq:factorblock} attains
$\FactorAutoExtremeARI{}$ and restores the lead. That is the expected direction:
at high noise the effective loadings of a module's members are most
heterogeneous, so the equal-loading assumption is least tenable. At $\sigma=1.5$
the two variants are close ($\FactorAutoHardARI{}$ versus $\MemeticHardARI{}$).

Two further points belong with these numbers. First, the margin over the best
competitor is modest---$0.03$ ARI at $\sigma=1.5$---where the margin over
average-linkage hierarchical clustering, the comparison a reader of the earlier
GA-clustering literature would expect, is $0.14$. A benchmark restricted to
agglomerative clustering and k-means would have overstated the advantage roughly
fourfold, and we report the stronger comparison because it is the honest one.
Second, the ablation isolating the memetic local search (Fig~\ref{fig:ablation}) is noise-dependent in the
same way: at $\sigma=1.5$ disabling it costs $0.065$ ARI ($p=2\times10^{-4}$),
confirming that the hill-climb---which reassigns genes \emph{a posteriori}, a move
greedy agglomeration cannot make---is responsible for much of the margin, while at
$\sigma=2.0$ the difference is not significant. Local search is decisive where the
likelihood surface is rugged but still informative, not universally.

A separate observation concerns WGCNA in its default configuration, which
performs poorly throughout the sweep (ARI $\WgcnaDefHardARI{}$ at $\sigma=1.5$).
The cause is diagnostic rather than incidental: with a mean within-module
$\lvert r\rvert$ near $0.29$, \texttt{pickSoftThreshold} selects a power of $8$,
at which the soft-thresholded adjacency is numerically negligible and the dynamic
tree cut returns a single grey module. This is a property of scale-free
thresholding on weakly correlated data rather than a deficiency of the method's
core idea, and it is why we compare against the oracle-tuned arm throughout.

% Place tables after the first paragraph in which they are cited.
\begin{table}[!ht]
\centering
\caption{\textbf{Module recovery against contemporary module detectors}
(ARI, mean $\pm$ SD over \NSeeds{} datasets) on the five-module design across a
noise sweep, with the number of clusters fixed to the truth or bisected to it.}
\label{tab:baselines}
\begin{adjustwidth}{-1.25in}{0in}
\centering
\IfFileExists{results/baselines_table.tex}{\input{results/baselines_table.tex}}{%
\begin{tabular}{@{}l@{}}\hline
\textit{Run experiments/benchmark\_baselines.py to populate this table.}\\\hline
\end{tabular}}
\end{adjustwidth}
\begin{flushleft}
Resolution- and inflation-based methods (Leiden, Louvain, MCL) do not accept a
cluster count; their parameters are bisected until they return the true $K$.
WGCNA selects its own module count, and is reported both at its default settings
and at the best of a power $\times$ minimum-module-size grid scored against
ground truth (Methods).
\end{flushleft}
\end{table}

\begin{figure}[!h]
\previewfig{baselines_vs_noise.pdf}
\caption{\textbf{Module recovery versus noise, against the full baseline set.}
Adjusted Rand index (mean, bands $\pm$ SD over \NSeeds{} datasets) at the true
number of clusters. \tool{} leads at moderate noise and has the lowest variance,
but the community-detection and network methods lie much closer to it than
agglomerative clustering does, and at $\sigma=2.0$ oracle-tuned WGCNA overtakes the constrained objective, though not the factor variant.}
\label{fig:perfnoise}
\end{figure}

\begin{figure}[!h]
\previewfig{ablation.pdf}
\caption{\textbf{The memetic local search helps at moderate noise, not
universally.} Recovery with and without the hill-climb across the noise sweep.
The refinement that reassigns elements after seeding accounts for much of
\tool{}'s margin at $\sigma\le1.5$; at $\sigma=2.0$ the difference is not
significant.}
\label{fig:ablation}
\end{figure}





\subsection*{When the model is right, and when it is not}
The benchmark above, like the synthetic designs in most module-detection papers,
draws its data from the same generative model the objective assumes. Success
there establishes that the optimizer reaches the likelihood optimum; it cannot
establish that the optimum is the biologically correct answer. We therefore
re-ran the comparison on the five processes of Methods, each violating a
different assumption (Table~\ref{tab:misspec}, Fig~\ref{fig:misspec}).

Across the \MisspecTotal{} processes at the discriminating noise level the
constrained objective leads on \MisspecWins{} and fails on two, and the pattern of
where it fails is the substantive result. Where modules are driven by a single
dominant signal it is competitive or better---its own one-factor design
$\MisspecOneFacGCM{}$, overlapping modules $\MisspecOverlapGCM{}$,
negative-binomial counts $\MisspecCountsGCM{}$, network-derived correlation
$\MisspecNetworkGCM{}$. The counts result is worth emphasizing: replacing the
Gaussian with a gamma--Poisson mixture and adding a mean--variance trend does not
degrade it, so \emph{distributional} misspecification is not what hurts this
model.

\emph{Structural} misspecification is. On multi-factor modules---where a module
spans a three-dimensional subspace rather than a single direction---it recovers
ARI $\MisspecMultiFacGCM{}$ where the best competitor reaches
$\MisspecMultiFacBest{}$, and on hub-and-spoke modules, whose members carry
systematically unequal loadings, $\MisspecHubGCM{}$ against $\MisspecHubBest{}$.
Both are exactly the regimes in which the equal-magnitude, single-factor
constraint is false.

That these are failures of the model rather than of the search can be shown
directly. On multi-factor data we evaluated Eq~\eqref{eq:loglik} at the planted
truth and at the partition \tool{} returns, and found the returned partition
scores \emph{higher} in \DiagLLMultiWins{} replicates (for example $\DiagLLExampleGCM{}$ against $\DiagLLExampleTruth{}$ at $\sigma=1.0$), while its mean ARI against the truth is only $\DiagLLMultiARI{}$; on hub-and-spoke data the same holds in \DiagLLHubWins{} replicates. The optimizer is
therefore working correctly and reaching a better optimum of the stated
objective; it is the objective whose optimum is in the wrong place. A likelihood
that outscores the ground truth is, in this sense, a sharper diagnostic of
misspecification than any accuracy metric, because it separates the two failure
modes that an accuracy drop alone conflates.

\begin{sidewaystable}
\centering
\caption{\textbf{Recovery under model misspecification} (ARI, mean $\pm$ SD over
\MisspecSeeds{} datasets per process, $\sigma=1.5$). Best method per process in
bold.}
\label{tab:misspec}
\IfFileExists{results/misspec_table_n15.tex}{\input{results/misspec_table_n15.tex}}{%
\begin{tabular}{@{}l@{}}\hline
\textit{Run experiments/misspecification.py to populate this table.}\\\hline
\end{tabular}}
\begin{flushleft}
The first column is \tool{}'s own generative model; the remaining five each
violate one of its assumptions (Methods). Every method receives the number of
modules planted in that dataset.
\end{flushleft}
\end{sidewaystable}

\begin{figure}[!h]
\previewfig{misspecification_n15.pdf}
\caption{\textbf{The advantage tracks whether modules carry a single dominant
signal.} Recovery by data-generating process. \tool{} leads where a module is one
direction in expression space---including when the data are counts rather than
Gaussian---and falls behind where a module spans several directions
(multi-factor) or where members carry systematically unequal loadings (hub).}
\label{fig:misspec}
\end{figure}

\subsection*{Freeing the loadings removes the failure on synthetic data}
The diagnosis predicts where relaxing the constraint should help. If the
constrained model fails because it cannot represent unequal or multi-directional
loadings, the free-loading model of Eq~\eqref{eq:factorblock} should recover the
loss, and its advantage over the constrained form should be concentrated on
exactly the two processes where Eq~\eqref{eq:loglik} fails---rather than being a
uniform improvement, which would suggest we had simply added capacity.

It does, and by more than repair. With the rank selected from the data by BIC (Methods), multi-factor recovery rises from
$\MisspecMultiFacGCM{}$ to $\MisspecMultiFacFactor{}$ and hub-and-spoke recovery
from $\MisspecHubGCM{}$ to $\MisspecHubFactor{}$. Both figures exceed the best
competing method on those processes ($\MisspecMultiFacBest{}$ for
\MisspecMultiFacBestName{} and $\MisspecHubBest{}$ for
\MisspecHubBestName{} respectively), so relaxing the exchangeability assumption
converts the framework's two clear failures into leads. On the four processes
where the exchangeable objective already led, the factor objective is equal or
slightly better---one-factor $\MisspecOneFacFactor{}$, overlapping
$\MisspecOverlapFactor{}$, counts $\MisspecCountsFactor{}$, network
$\MisspecNetworkFactor{}$---so it leads on all \MisspecTotal{} processes at this
noise level, where the exchangeable objective leads on \MisspecWins{}
(Table~\ref{tab:misspec}). The same pattern appears in the noise sweep, where the
factor objective is the only variant to beat oracle-tuned WGCNA at $\sigma=2.0$.

Rank selection recovers the generating structure without supervision: BIC chooses
$q=\MisspecMultiFacRank{}$ on the three-factor process and
$q=\MisspecOneFacRank{}$ on the one-factor, hub, overlapping, counts and network processes; in a separate experiment the automatically selected rank matched the recovery of the best rank chosen with knowledge of the truth in \DiagAutoMatches{} fits (mean shortfall $\DiagAutoGap{}$ ARI). The extra flexibility therefore costs computation but
not accuracy.

The rank must be selected rather than fixed. Setting $q$ above the number of
factors actually present is worse than not using the branch at all: at $q=3$ on
the one-factor design recovery falls from $\DiagOneFacLoglikARI{}$ to $\DiagOneFacQthreeARI{}$, because the surplus
factors fit noise directions within each block. This is why the size-scaled
parameter count of Eq~\eqref{eq:factorblock} matters---it is what makes BIC
comparable across ranks and therefore makes automatic selection possible. We
recommend BIC over AIC here, since at the planted partition AIC selected one rank too many on the three-factor process (correct in \DiagRankAicMultiFac{} datasets, against \DiagRankBicMultiFac{} for BIC).

The factor branch is offered alongside the exchangeable objective rather than as
a replacement. On synthetic designs it costs roughly three to four times more per
run, multiplied by the number of candidate ranks; on real compendia the gap is far
larger, between \SaelensCostRatioMin{} and \SaelensCostRatioMax{} times (Results), because real correlation
structure yields more modules and a higher selected rank, and on data that genuinely follow the one-factor
model it offers nothing over Eq~\eqref{eq:loglik}. Its value is that a user who
has specific reason to expect module members to respond at different strengths, or to several regulators, can test that expectation; on the real compendia analysed below we found no such case.

\subsection*{Real expression data with curated modules}
Everything above is measured on data we generated. The misspecification sweep
mitigates that---those processes deliberately break our own model---but they
remain simulations, and a module detector should be judged on whether it recovers
modules that biologists have independently curated. We therefore evaluated on an established module-detection benchmark~\cite{Saelens2018}, which pairs real expression
compendia with module definitions derived from regulatory networks:
RegulonDB regulons for \emph{E. coli}, MacIsaac and Sisima regulons for yeast, and
regulatory circuits for human (Methods).

Ground truth here is not a partition. Regulons overlap and do not cover every
gene, so the adjusted Rand index is inapplicable and we use the benchmark's own
scores: \emph{recovery}, the mean over known modules of the best Jaccard against
any predicted module; \emph{relevance}, the mean over predicted modules of the
best Jaccard against any known module; and their harmonic mean, F1. Recovery
alone is maximized by predicting one enormous module and relevance alone by
predicting many tiny ones, so F1 is the score to read.

Over \SaelensNDatasets{} datasets spanning three organisms
(Table~\ref{tab:saelens}, Fig~\ref{fig:saelens}), \tool{} driven by the
constrained criterion attains the best mean F1 of any method,
$\SaelensLoglikScore{}$, and wins \SaelensLoglikWins{} of the \SaelensNDatasets{}
datasets outright. Against the graph and agglomerative methods the result is
clean: it beats Leiden and average-linkage hierarchical clustering on every
dataset ($\SaelensLeidenBeaten{}$, sign test $p=\SaelensLeidenP{}$; mean F1
$\SaelensLeidenScore{}$ and $\SaelensHierScore{}$ respectively) and spectral clustering
on $\SaelensSpectralBeaten{}$. Against WGCNA it does not separate: \tool{} leads
on $\SaelensWgcnaBeaten{}$ datasets with a mean F1 of $\SaelensLoglikScore{}$ against
$\SaelensWgcnaScore{}$, a difference that a sign test over six datasets cannot
distinguish from chance ($p=\SaelensWgcnaP{}$). The defensible claim is therefore
that \tool{} attains the best mean recovery of curated modules among the methods
tested, is clearly ahead of the community-detection and agglomerative approaches,
and is on a par with WGCNA.

The free-loading factor model, which led every synthetic process, does not
transfer. It is worse than the constrained criterion on \SaelensFactorBeaten{} real datasets (sign test $p=\SaelensFactorP{}$; mean F1
$\SaelensFactorScore{}$ against $\SaelensLoglikScore{}$), never wins one, and also
trails WGCNA. Taken with the breast-cancer analysis below, where it likewise
fails to improve on the constrained form, this is consistent evidence across
seven independent real datasets. We take it at face value: the structures that
the factor model is built for---modules spanning several latent directions, or
whose members respond at very different strengths---are ones we can construct in
simulation and detect there, but they do not appear to characterize these
compendia at the scale analysed, so the additional free parameters add variance
without capturing signal. This is a statement about real co-expression data, not
about the optimizer, and it is the reason we recommend Eq~\eqref{eq:loglik} as
the default despite the synthetic results.

Two caveats bound these numbers. We analysed the \SaelensNTop{} most variable
genes per dataset rather than whole compendia, because the factor procedure costs
between \SaelensCostRatioMin{} and \SaelensCostRatioMax{} times more than the constrained criterion at this
scale and whole-compendium runs with it were not tractable; absolute scores are
therefore not comparable with those published by Saelens et al., though the
between-method comparison, which is what we claim, is unaffected. And absolute F1
is low throughout---from $\SaelensFMin{}$ on human to $\SaelensFMax{}$ on \emph{E. coli}---which
reflects the genuine difficulty of recovering regulons from expression alone
rather than a defect of any particular method; the human compendia are the
hardest for every method tested.

\begin{table}[!ht]
\centering
\caption{\textbf{Recovery of curated modules on the Saelens et al. benchmark}
(F1 of recovery and relevance; higher is better). Best method per dataset in
bold. All methods see the same \SaelensNTop{} most variable genes per dataset.}
\label{tab:saelens}
\begin{adjustwidth}{-1.25in}{0in}
\centering
\IfFileExists{results/saelens_table.tex}{\begin{tabular}{@{}lr@{}}
\hline
Method & ecoli\_colombos \\
\hline
\tool{} (\texttt{loglik\_factor}, two-stage) & 0.148 \\
\tool{} (\texttt{loglik\_aic}) & \textbf{0.265} \\
WGCNA & 0.029 \\
Leiden & 0.171 \\
spectral & 0.237 \\
hierarchical (average) & 0.136 \\
\hline
\end{tabular}
}{%
\begin{tabular}{@{}l@{}}\hline
\textit{Run experiments/saelens\_benchmark.py to populate this table.}\\\hline
\end{tabular}}
\end{adjustwidth}
\end{table}

\begin{figure}[!h]
\previewfig{saelens.pdf}
\caption{\textbf{Module recovery on real expression compendia with curated
ground truth.} F1 of recovery and relevance by dataset and method. The
constrained correlation criterion attains the best mean score and is clearly
ahead of Leiden, spectral and agglomerative clustering; it is not distinguishable
from WGCNA. The free-loading factor model, which leads every synthetic process,
is worse than the constrained form on all six real datasets.}
\label{fig:saelens}
\end{figure}

\subsection*{Selecting the number of clusters without knowing it}
When only \texttt{loglik} is optimized under a loose $g_{\max}$, the correlation
objective can over-segment, since Eq~\eqref{eq:loglik} carries no parsimony term.
The effect, however, is mild once the likelihood is optimized well: on
Synthetic-Easy the optimum already coincides with the true partition
($k=\EasyLooseK{}$), and on Synthetic-Hard a loose bound yields only
$k=\HardLooseK{}$ (true 5; Table~\ref{tab:benchmark}), one module above the truth.
The penalized correlation criteria (Eq~\eqref{eq:loglikic}) keep selection near
the truth within a single objective: on Synthetic-Easy \texttt{loglik\_aic}
returns $k=\EasyAicK{}$ at ARI $=\EasyAicARI{}$, and on this Synthetic-Hard
dataset $k=\HardAicK{}$ at ARI $=\HardAicARI{}$. The decisive test is replication, and it is
also where this route shows its limit. Over \NSeeds{} datasets
\texttt{loglik\_aic} selects $\MemeticSelK{}$ modules on average at ARI
$\MemeticSelARI{}$, comfortably better than hierarchical clustering whose cut
minimizes the same criterion ($\HierSelK{}$ modules, ARI $\HierSelARI{}$;
Table~\ref{tab:modelsel}), so global optimization does help when the criterion is
held fixed. Against the community-detection methods, however, the picture is less
favourable: free-running Leiden at its default resolution recovers
$\LeidenSelK{}$ modules at ARI $\LeidenSelARI{}$, statistically indistinguishable
from \texttt{loglik\_aic} in ARI but substantially closer to the true count
($\lvert k-5\rvert$ of $0.27$ against $1.73$, $p=2\times10^{-5}$). A resolution
default with no explicit model therefore identifies the number of modules more
reliably than our penalized correlation criterion does, and we do not claim
otherwise.

The reason is instructive, and it also rules out the obvious remedy. The penalty
in Eq~\eqref{eq:loglikic} charges a fixed cost per module, because the
exchangeable block model has one parameter per module irrespective of size;
splitting a large module is therefore cheap relative to the likelihood it buys,
which produces the mild over-segmentation observed. It is natural to expect the
factor block model of Eq~\eqref{eq:factorblock} to do better, since its parameter
count grows with module size. It does not---it is substantially worse, selecting
the largest number of modules permitted in every replicate
(Table~\ref{tab:modelsel}). The cause is a property of the parameter count that
is easy to miss. Summing $p_m = n_m q - q(q-1)/2 + 1$ over $k$ modules whose
sizes total $n$ gives
\begin{equation}
\textstyle\sum_m p_m \;=\; nq \;-\; k\,q(q-1)/2 \;+\; k ,
\label{eq:factorparams}
\end{equation}
which is increasing in $k$ only for $q=1$; it is constant in $k$ at $q=2$ and
\emph{decreasing} for $q\ge3$. The per-block correction $q(q-1)/2$ for rotational
invariance is correct for each block in isolation, but accumulating it across
blocks means that at higher ranks a partition is charged \emph{less} for having
more modules. The criterion therefore cannot penalize partition complexity, and
the search runs to whatever bound $g_{\max}$ imposes. We report this as a
limitation of the factor branch rather than repair it here: the underlying issue
is that BIC accounts for the parameters of a partition but not for the
combinatorial multiplicity of the partition itself, which the fixed per-module
term of Eq~\eqref{eq:loglikic} penalizes only incidentally. \textbf{The factor
objective should therefore be used with $K$ known or externally constrained, and
Eq~\eqref{eq:loglikic} retained when the module count must be inferred.} Users
whose priority is recovering the module count itself, rather than optimizing a
stated correlation criterion, should treat a resolution-based method as a strong
and much cheaper baseline.

This division of labour suggests a practical recipe, since the two objectives
fail in complementary ways: use the constrained criterion, whose penalty does
grow with the module count, to choose $K$, then fit the factor model with $K$
held at that value. Over \TwoStageSeeds{} datasets of the five-module design at
$\sigma=1.5$ this two-stage procedure reaches ARI $\TwoStageARI{}$, against
$\TwoStageAicARI{}$ for \texttt{loglik\_aic} alone and $\TwoStageLooseARI{}$ for
the factor model run with a loose $g_{\max}$, and approaches the
$\TwoStageOracleARI{}$ obtained when the true $K$ is supplied. The residual gap is
inherited from the first stage, which selects $\TwoStageK{}$ modules on average
against a truth of five; a better $K$-selection procedure---a resolution-based
method, for instance---would close it further. We recommend this two-stage
workflow whenever the module count is unknown. The
fit--parsimony trade-off is also visible by scanning $g_{\max}$
(Fig~\ref{fig:modelsel}): \texttt{loglik} rises and plateaus with $k$ while the penalized criteria fall steeply up to the true number of modules and then flatten, reaching their minimum one module above it.

\begin{figure}[!h]
\caption{\textbf{Model selection on Synthetic-Hard.}
The raw correlation log-likelihood (left axis) increases monotonically as
clusters are added---the over-segmentation tendency---while the penalized
criteria (right axis, lower is better) fall steeply up to the true number of modules and then flatten, reaching their minimum one module above it (vertical line). The horizontal axis is the number of clusters the optimizer returned as $g_{\max}$ was raised; it saturates at six, so larger values do not appear.}
\previewfig{model_selection.pdf}
\label{fig:modelsel}
\end{figure}

For users who prefer to inspect an explicit trade-off rather than commit to a
single penalized criterion, the multi-objective mode optimizes \texttt{loglik}
against a parsimony target under NSGA-II and returns the full Pareto front
(Fig~\ref{fig:pareto}). We note that pairing \texttt{loglik} with the
\emph{geometric} BIC is less effective on these data, because BIC's
Euclidean-Gaussian assumption conflicts with the absolute-correlation objective
when modules contain anti-correlated members---which is exactly why the penalized
\emph{correlation} criteria are preferable for module selection.

\begin{figure}[!h]
\caption{\textbf{Fit--parsimony Pareto front on Synthetic-Hard.}
Each point is a non-dominated partition from a single NSGA-II run optimizing the
correlation log-likelihood (maximized) against the geometric BIC (minimized);
colour encodes the number of clusters. The multi-objective mode exposes the
trade-off explicitly for interactive exploration.}
\previewfig{pareto.pdf}
\label{fig:pareto}
\end{figure}

\begin{table}[!ht]
\centering
\caption{\textbf{Model selection with $k$ unknown} (mean $\pm$ SD over \NSeeds{} datasets, five-module design at $\sigma=1.5$, true $k=5$). Every method runs free: no ground-truth $k$ and no oracle tuning. The \tool{} and hierarchical rows use the same penalized correlation criterion and differ only in optimizer.}
\label{tab:modelsel}
\IfFileExists{results/baselines_modelsel.tex}{\input{results/baselines_modelsel.tex}}{%
\begin{tabular}{@{}l@{}}\hline
\textit{Run experiments/benchmark\_baselines.py to populate this table.}\\\hline
\end{tabular}}
\end{table}

\subsection*{Robustness to unstructured genes}
Real gene sets contain many genes that belong to no module, and this is where
\tool{}'s margin is largest and least ambiguous. When the five-module design is
embedded among unstructured ``noise'' genes and the methods must both find the
modules and avoid being derailed by the noise, \tool{} retains ARI
$\NoiseGeneMemeticARI{}$, against $\NoiseGeneHierARI{}$ for hierarchical
clustering under the same penalized criterion, $\NoiseGeneLeidenARI{}$ for
free-running Leiden and $\NoiseGeneWgcnaARI{}$ for WGCNA
(\NSeeds{} datasets; every paired difference $p<10^{-4}$). Unlike the noise-sweep
margin, which is a few hundredths of ARI, these differences are an order of magnitude larger.

The mechanism is the same likelihood that fails on multi-factor data, working in
its favour here. Because a module's contribution to Eq~\eqref{eq:loglik} depends
on its mean within-module $\lvert r\rvert$, admitting an incoherent gene lowers
that mean and is penalized; the optimizer consequently isolates unstructured
genes into their own low-coherence groups rather than distributing them among
real modules. Greedy linkage, by contrast, must place every gene at some merge
height and forces them into the nearest module, and WGCNA's grey module absorbs
them only when the tree cut happens to separate them. For the common practical
situation in which a gene set has been filtered by variance rather than by
membership in any known programme, this behaviour is arguably more valuable than
a small gain in recovery on clean data.

\subsection*{Scaling to realistic gene counts}
Co-expression studies routinely cluster thousands of genes, so the practical
question is not only whether a global optimizer is more accurate but whether it
is affordable at that size. We measured runtime and recovery on the modular
design at $200$ to $\ScaleMaxGenes{}$ genes (Table~\ref{tab:scaling}).

Recovery is unaffected by size (ARI $\ge0.996$ throughout). Runtime at
$\ScaleMaxGenes{}$ genes is $\ScaleMaxSeconds{}$~s per dataset, and over the
$1000$--$\ScaleMaxGenes{}$ range it grows as $n^{\ScaleExponent{}}$---close to
linear, and well below the $O(n^2)$ cost of forming the correlation matrix that
every correlation-based method must pay. \tool{} is accordingly two to three
times slower than Leiden and WGCNA at $\ScaleMaxGenes{}$ genes rather than orders
of magnitude slower.

These timings are for the constrained objective of Eq~\eqref{eq:loglik}, and the
factor objective is considerably more expensive---not by the constant factor of
three to four that its per-evaluation cost implies, but by between \SaelensCostRatioMin{} and \SaelensCostRatioMax{} times on real data. On the curated compendia the constrained criterion fitted $\SaelensNTop{}$ genes in $\SaelensLoglikSecMin{}$ to $\SaelensLoglikSecMax{}$~s where the factor stage took $\SaelensFactorSecMin{}$ to $\SaelensFactorSecMax{}$~s, because the local search must re-solve an eigendecomposition
per affected block and real correlation structure yields both more modules and a
higher selected rank than the synthetic designs do. Together with the recovery
results this is why we recommend the constrained criterion as the default: on
real data it is both better and far cheaper.

We measured only to $\ScaleMaxGenes{}$ genes, and we do not extrapolate: the
$O(n^2)$ correlation matrix becomes the binding constraint before the optimizer
does, and at $16{,}000$ genes it alone occupies roughly two gigabytes in double
precision. Genome-wide application therefore still benefits from the variance
pre-filtering used here, though for the constrained objective the ceiling is set
by memory for the correlation matrix rather than by the search.

\begin{table}[!ht]
\centering
\caption{\textbf{Runtime scaling} (seconds per dataset, mean of three
replicates). Recovery is unaffected by size across this range.}
\label{tab:scaling}
\IfFileExists{results/scaling_table.tex}{\input{results/scaling_table.tex}}{%
\begin{tabular}{@{}l@{}}\hline
\textit{Run experiments/scaling.py to populate this table.}\\\hline
\end{tabular}}
\end{table}

\subsection*{Optimizer validation on geometric data: Iris}
We include Iris not as a recovery benchmark---it is geometric rather than
correlational, outside the correlation objective's domain---but to verify that
the GA is a faithful, competent optimizer of \emph{any} geometric index it is
given. It is: driven by \texttt{silhouette} it returns a partition with
silhouette $\IrisGcmrkSil{}$, matching exactly the optimum found by k-means at
$k=2$ ($\IrisKmTwoSil{}$), so the optimizer reaches the index's best value rather
than a local one.

That this partition has $k=\IrisSilK{}$ rather than three clusters is a property
of the \emph{index}, not the optimizer: on Iris the three-species labelling has
silhouette only $\IrisTrueSil{}$, \emph{below} both the two-cluster split
($\IrisKmTwoSil{}$) and the best three-cluster k-means partition
($\IrisKmThreeSil{}$), because \emph{versicolor} and \emph{virginica} overlap so
heavily that merging them is geometrically cleaner (Fig~\ref{fig:iris}). No
silhouette-maximizing method can return the three species here; \tool{} correctly
returns what the criterion prefers (ARI $=\IrisSilARI{}$ against the species,
matching k-means at the same $k$). The lesson is that the choice of criterion,
not the optimizer, determines the answer, which is precisely the control \tool{}
gives the user, and the reason the correlation objective matters for data where a
geometric index would be the wrong question.

\begin{figure}[!h]
\caption{\textbf{Silhouette profile of the Iris clustering recovered by \tool{}.}
Per-sample silhouette coefficients grouped by recovered cluster; the dashed line
marks the mean. One cluster is cleanly separated (uniformly high silhouette)
while the other two overlap, the expected structure for the Iris species.}
\previewfig{iris_silhouette.pdf}
\label{fig:iris}
\end{figure}

\subsection*{Empirical co-expression modules: breast cancer}
We applied every method to the \EmpCompNTop{} most variable genes of
GSE183947~\cite{Zhang2021}, a breast-cancer RNA-seq dataset of \EmpSamples{}
samples (30 tumours and 30 patient-matched adjacent-normal tissues), and scored
the resulting partitions by four criteria, none of which uses the tumour label to
fit anything (Table~\ref{tab:empcomp}, Fig~\ref{fig:empcomp}).

The first criterion is the one the earlier literature, and an earlier version of
this manuscript, relied on: whether module eigengenes separate tumour from normal
tissue. It does not discriminate between methods. Every method produces modules
whose eigengenes differ significantly between tissue types
(Benjamini--Hochberg $q<0.05$; Fig~\ref{fig:empirical}), which is unsurprising given that the
tumour/normal contrast dominates variation in this dataset; almost any partition
of its most variable genes will recover some of it. We report this criterion for
continuity with the literature while noting that it is close to unfalsifiable and
should not be used to argue that one module detector is better than another.

The remaining three criteria are more informative, and they do not favour any
method decisively. Modules from \tool{} are as coherent as those from the
alternatives (median within-module $\lvert r\rvert$ $\EmpCompLoglikCoh{}$ for the
constrained objective and $\EmpCompFactorCoh{}$ for the factor model, against
$\EmpCompWgcnaCoh{}$ for WGCNA and $\EmpCompLeidenCoh{}$ for Leiden), though this
comparison slightly favours \tool{} by construction, since coherence is close to
what the constrained objective maximizes. On functional enrichment---the
criterion that asks whether a module corresponds to a recognizable biological
programme---\tool{} enriches $\EmpCompLoglikEnr{}$ of $\EmpCompLoglikTested{}$
tested modules against $\EmpCompLeidenEnr{}$ of $\EmpCompLeidenTested{}$ for
Leiden and $\EmpCompWgcnaEnr{}$ of $\EmpCompWgcnaTested{}$ for WGCNA. The
proportions favour \tool{}, but the counts are small enough that we draw no
conclusion from them.

Reproducibility is the criterion on which \tool{} does not lead. Splitting the
samples in half, clustering each half independently and comparing the two
partitions gives an adjusted Rand index of $\EmpCompLeidenSplit{}$ for Leiden
against $\EmpCompLoglikSplit{}$ and $\EmpCompFactorSplit{}$ for the two \tool{}
variants. That comparison is confounded, because \tool{} selects
$\EmpCompLoglikK{}$ modules here where Leiden selects $\EmpCompLeidenK{}$, and
agreement measured on finer partitions is systematically lower; but the confound
does not favour us, and the honest summary is that \tool{}'s modules are not more
stable than a resolution-based method's on this dataset. That \tool{} selects
roughly three times as many modules as Leiden is itself consistent with the
over-segmentation of the penalized correlation criterion documented above.

Two further observations deserve recording. First, the factor model does
\emph{not} outperform the constrained objective here, on any of the four
criteria, despite leading on every synthetic process---the same pattern seen on
all \SaelensNDatasets{} datasets of the curated benchmark above. Across the seven
real datasets analysed in this paper the free-loading model never once improves on
the constrained one, which we take as evidence that modules in real compendia at
this scale behave as single-signal structures. Second,
the absolute reproducibility of \emph{every} method is low
($\text{ARI} \le \EmpCompLeidenSplit{}$), which is a statement about the dataset
rather than the methods: with 30 samples per half, module structure among
\EmpCompNTop{} highly variable genes is simply not sharply determined. A
convincing demonstration that any of these methods recovers more biology than the
others would need a substantially larger cohort, and we regard that as the main
open item for this line of work.

\begin{table}[!ht]
\centering
\caption{\textbf{Comparative analysis of GSE183947.} All methods applied to the
same \EmpCompNTop{} most variable genes. ``assigned'' is the fraction of genes
placed in a module (WGCNA leaves some in its grey module); ``GO-enriched
modules'' counts modules with at least one \EmpCompLibrary{} term at adjusted
$p<0.05$ against the analysed gene set as background; ``split-half ARI'' compares
partitions fitted independently to two halves of the samples.}
\label{tab:empcomp}
\begin{adjustwidth}{-1.25in}{0in}
\centering
\IfFileExists{results/empirical_comparison.tex}{\input{results/empirical_comparison.tex}}{%
\begin{tabular}{@{}l@{}}\hline
\textit{Run experiments/empirical\_comparison.py to populate this table.}\\\hline
\end{tabular}}
\end{adjustwidth}
\end{table}

\begin{figure}[!h]
\previewfig{empirical_comparison.pdf}
\caption{\textbf{No method dominates on real data.} Median within-module
correlation, fraction of modules with significant GO enrichment, and split-half
reproducibility for each method on GSE183947. \tool{} (dark bars) leads on
coherence and enrichment fraction and trails on reproducibility; none of the
differences is large relative to what a single dataset of this size can resolve.}
\label{fig:empcomp}
\end{figure}

\begin{figure}[!h]
\caption{\textbf{Module eigengenes separate tumour from normal tissue.}
Eigengene distributions (leading principal component of each module's expression
across the \EmpSamples{} samples) for the \tool{} modules most associated with
phenotype, split by tumour (T) and patient-matched normal (N) tissue; asterisks
mark modules with a significant contrast (Benjamini--Hochberg $q<0.05$). Every
method tested produces such modules, so this figure illustrates the output rather
than discriminating between methods.}
\previewfig{empirical_eigengenes.pdf}
\label{fig:empirical}
\end{figure}

\subsection*{Visualizing cluster structure in reduced dimensions}
To provide visual confirmation that the recovered modules occupy distinct regions
of the space in which the objective operates, we embed each clustering result
into two dimensions using PCA and correlation MDS
($1 - \lvert\mathbf{R}\rvert$; see Methods).  On Synthetic-Easy the three
recovered modules form well-separated clouds in both embeddings
(Fig~\ref{fig:dimred_synth}), consistent with the perfect ARI.  On
Synthetic-Hard, where noise and unequal module sizes degrade separation, PCA
shows considerable overlap because co-expressed but anti-correlated genes are
geometrically distant; the correlation MDS, which measures distance in the
same absolute-correlation space as the objective, reveals clearer module
boundaries.  For the empirical breast-cancer data, the correlation-MDS
embedding of the \EmpKmodules{} modules shows that even on real gene expression
data the recovered modules occupy coherent, distinguishable regions of
correlation space (Fig~\ref{fig:dimred_emp}).  These visualizations are produced
automatically by the software and are available to users for any dataset.

\begin{figure}[!h]
\caption{\textbf{Dimensionality-reduction visualization of synthetic
clustering.} Recovered modules projected into two dimensions by PCA (left) and
correlation MDS (right) on Synthetic-Easy.  Diamonds mark cluster centroids.
The correlation-MDS embedding uses the distance space native to the
correlation objective and shows clean separation.}
\previewfig{dimred_synthetic_easy.pdf}
\label{fig:dimred_synth}
\end{figure}

\begin{figure}[!h]
\caption{\textbf{Dimensionality-reduction visualization of empirical modules.}
The \EmpKmodules{} co-expression modules recovered from GSE183947 projected into
two dimensions by PCA (left) and correlation MDS (right).  Module boundaries are
more distinct in the correlation distance space, confirming that the objective
captures structure beyond what Euclidean geometry reveals.}
\previewfig{dimred_empirical.pdf}
\label{fig:dimred_emp}
\end{figure}

\section*{Discussion}
\tool{} reframes co-expression module detection as maximum-likelihood inference
under an explicit generative model, and the value of making that model explicit
is that it can be tested rather than assumed. The correlation criterion the field
already uses is not a heuristic: it is, up to a constant and the sample-size
factor, the profile log-likelihood of a block-diagonal one-factor Gaussian in
which each module is driven by a single regulator with equal-magnitude $\pm$
loadings \cite{bartholomew2011} (\nameref{AppendixA}). Writing it that way makes
two constraints visible---one factor, equal magnitudes---and both can be relaxed
within the same family. Whether to relax them is then an empirical question
rather than a matter of taste, and our answer is split. On synthetic data the
free-loading model of Eq~\eqref{eq:factorblock} is at least as good as the
constrained form everywhere and strictly better where the constraints bind. On
real expression compendia with curated modules it is worse on every dataset we
tested, never wins one, and costs \SaelensCostRatioMin{} to \SaelensCostRatioMax{} times more. We therefore
recommend the constrained absolute-correlation criterion as the default, and the
factor model where there is specific reason to think module members respond at
very different strengths or to several regulators, which is a situation we can construct
and detect in simulation but did not find in seven real datasets.

We think the negative transfer is itself worth reporting. The obvious reading is
that our generators overstate how often real modules are multi-directional: a
regulon at the scale of a few hundred variable genes appears to behave much like
a single dominant signal, so the extra loadings fit noise rather than structure.
That is a claim about co-expression data, not about the optimizer, and it is
testable by anyone who finds a compendium where the factor model does win. It
also cautions against a common inference in methods papers, that a model which
dominates on simulations designed to exercise it will dominate in application.

The likelihood footing also explains why pairing
either objective with a \emph{geometric} information criterion fails: those assume
Euclidean-Gaussian clusters, which a correlation block model contradicts whenever
modules contain anti-correlated genes.

Given the model, the first contribution that matters in practice is optimization
quality. Greedy agglomerative clustering on the correlation distance (the core
of widely used co-expression pipelines) is a strong seed but a poor optimizer:
once it merges two genes it can never separate them. Our replicated benchmark
shows that hybridizing the GA with a memetic local search, which reassigns genes
after the fact, surpasses hierarchical clustering and every contemporary
alternative we tested at moderate noise. We are deliberate about the size of that
claim. Measured against agglomerative clustering alone the margin is $0.14$ ARI;
measured against oracle-tuned WGCNA, the strongest competitor, it is $0.03$, and
at $\sigma=2.0$ it overtakes the constrained objective, though not the factor variant.

The second contribution is a negative one that we think is more useful than
another favourable table. By generating data from processes the model does not
assume, we found two structural regimes---modules spanning several latent factors,
and modules whose members carry systematically unequal loadings---in which
maximizing the correlation likelihood moves away from the true partition. That
this is misspecification rather than search failure is directly demonstrable: the
optimizer returns partitions of higher likelihood than the ground truth. The
diagnostic generalizes beyond this method. Any model-based clustering procedure
can be interrogated the same way, by evaluating its own objective at a trusted
reference partition; an objective that outscores the truth is misspecified for
those data, and no amount of better search will help. Notably, distributional
misspecification (non-Gaussian counts with a mean--variance trend) did not
degrade performance, which suggests that practitioners should worry more about
whether their modules are single-signal than about whether their data are normal.

The third contribution follows from the second. Because the failure is
attributable to a specific assumption (exchangeability of module members) it
can be repaired by relaxing that assumption rather than abandoning the framework.
The rank-$q$ free-loading block model does not merely recover the lost accuracy:
with the rank selected from the data it leads every synthetic process we tested, including the two the exchangeable objective fails on, and it is the only variant that beats oracle-tuned WGCNA at the highest noise level---though, as the curated benchmark showed, none of this carries over to real compendia. Because its parameter count
scales with module size, its information criteria are comparable across ranks, so
the rank is selected without supervision. It does not, however, help with
selecting the number of modules: as Eq~\eqref{eq:factorparams} shows, its
parameter count stops increasing with the module count at $q\ge2$, so the
criterion that selects the rank so successfully cannot also select $K$. The two
selection problems are not solved by the same penalty, and we recommend the
factor objective only where $K$ is known or bounded.

Limits remain, and two deserve emphasis. First, our penalized correlation
criterion is not the best available way to choose the number of modules: a
resolution-based community detector with no explicit model recovers the module
count more reliably. The contribution here is a principled criterion within a
likelihood framework, not a state-of-the-art model-selection procedure. Second,
scaling is bounded by the $O(n^2)$ correlation matrix rather than by the search,
so genome-wide application still benefits from pre-filtering. Natural extensions
include shrinkage-regularized within-module covariance, correlation-aware
crossover operators, and sketched or approximate correlation for very large $n$.

Methodologically, \tool{} sits in the tradition of GA-based and multi-objective
clustering~\cite{Maulik2000,Handl2007} but is distinguished by a co-expression
objective with an explicit likelihood, a memetic optimizer that measurably
outperforms the greedy clustering such pipelines rely on, and---unusually---a
characterization of the conditions under which its own objective should not be
used. We note that multi-objective search is not a remedy for the
misspecification we identify: a Pareto front between two objectives helps only
when at least one of them ranks the truth highly, and graph modularity, the natural second objective, is itself misspecified on the same processes: evaluated at the planted partition it ranks the truth above the partitions \tool{} and Leiden return in only \DiagModTruthTopMultiFac{} multi-factor and \DiagModTruthTopHub{} hub-and-spoke datasets, against \DiagModTruthTopOneFac{} on the one-factor design (Newman modularity on the soft-thresholded graph, which is not the exact objective Leiden optimizes). Search strategy cannot repair a wrongly located optimum. As a
practical complement, the built-in dimensionality-reduction tools (PCA and
correlation MDS, with optional t-SNE and UMAP) let users inspect and communicate
the recovered structure, with the correlation-MDS embedding particularly
informative because it projects elements into the same distance space the
objective operates in.

\section*{Conclusion}
We introduced \tool{}, a dependency-light Python tool that detects co-expression
modules by maximizing the likelihood of an explicit block-diagonal Gaussian model
with a memetic genetic algorithm. Recognizing the field's usual
absolute-correlation criterion as the constrained member of a factor-model family
puts the modelling assumptions on the table, where they can be tested rather than
inherited. Tested, they largely hold: on the Saelens benchmark of real expression
compendia with curated modules, the constrained criterion attains the best mean
recovery of any method over \SaelensNDatasets{} datasets in three organisms,
beating Leiden and agglomerative clustering on every one and matching WGCNA; on
replicated synthetic designs it leads every contemporary detector at moderate
noise; and it isolates genes belonging to no module far better than anything else
we tested.

Making the model explicit also let us find where it breaks. We identify two
structural regimes in which the constrained objective is misspecified,
demonstrate that its optimizer there returns partitions scoring above the ground
truth, a diagnostic that separates a wrong model from a failed search, and that
applies to any model-based clustering method, and show that freeing the loadings
turns both failures into leads on those processes. We then report that this
repair does not transfer to real data, that the criterion which selects the
factor rank cannot select the number of modules, and that our advantage over
WGCNA on real compendia is not statistically distinguishable. We suggest that
reporting where a method's objective ceases to be appropriate, and where an
improvement fails to generalize, is more useful to practitioners than a benchmark
in which it always wins. The software, documentation, and full reproduction
scripts are openly available.

\section*{Appendices}

\paragraph*{Appendix A.}
\namedlabel{AppendixA}{Appendix A}
\textbf{Derivation of the correlation objective from the one-factor block model.}

\emph{Model.} Each of the $d$ samples (columns of $\mathbf{X}$) is an
independent draw $\mathbf{x}\sim N(\mathbf{0},\Sigma)$ over the $n$
row-standardized genes, with $\Sigma$ block-diagonal by the partition. Within a
module $s$ of size $n_s$, the equal-magnitude one-factor structure
$x_g=\sigma_g f_s+\varepsilon_g$, $\sigma_g\in\{\pm\lambda_s\}$, induces a
correlation matrix that, after absorbing the loading signs by a diagonal
$\pm1$ similarity (which changes neither eigenvalues nor determinant), is the
equicorrelation (compound-symmetry) matrix
$\Sigma_{\mathrm{CS}}(\rho_s)=(1-\rho_s)I+\rho_s\,\mathbf{1}\mathbf{1}^{\!\top}$
with $\rho_s=\lambda_s^2$. Its eigenvalues are $1+(n_s-1)\rho_s$ once and
$1-\rho_s$ with multiplicity $n_s-1$, so
$\log\det\Sigma_{\mathrm{CS}}(\rho_s)=\log[1+(n_s-1)\rho_s]+(n_s-1)\log(1-\rho_s)$.

\emph{Profile likelihood.} For $d$ i.i.d. samples with sample correlation
$\mathbf{S}_s$ the per-module Gaussian log-likelihood is
$\ell_s(\rho)=-\tfrac{d}{2}\big[\log\det\Sigma_{\mathrm{CS}}(\rho)+\operatorname{tr}(\Sigma_{\mathrm{CS}}(\rho)^{-1}\mathbf{S}_s)+n_s\log 2\pi\big]$.
Writing $P=\mathbf{u}^{\!\top}\mathbf{S}_s\mathbf{u}$ for the variance along the
aligned direction $\mathbf{u}=\mathbf{s}/\sqrt{n_s}$ (sign vector $\mathbf{s}$),
the two distinct eigenvalues of $\Sigma_{\mathrm{CS}}$ are matched at their ML
values $\widehat{1+(n_s-1)\rho}=P$ and $\widehat{1-\rho}=(n_s-P)/(n_s-1)$,
giving $\hat\rho_s=(P-1)/(n_s-1)$. With sign alignment
$P=1+(n_s-1)\bar r_s$ where $\bar r_s$ is the mean off-diagonal $\lvert R\rvert$,
hence $\hat\rho_s=\bar r_s=(C_s-n_s)/(n_s^2-n_s)$ (Eq~\eqref{eq:cs}). At the ML
estimate the trace term collapses, $\operatorname{tr}(\Sigma_{\mathrm{CS}}(\hat\rho_s)^{-1}\mathbf{S}_s)=n_s$,
so
$\ell_s=-\tfrac{d}{2}\log\det\Sigma_{\mathrm{CS}}(\hat\rho_s)-\tfrac{d}{2}n_s(1+\log 2\pi)$.

\emph{Objective.} Summing over modules, the second term is
$-\tfrac{d}{2}(1+\log 2\pi)\sum_s n_s$, a constant (total genes fixed). The
partition-dependent part is
$-\tfrac{d}{2}\sum_s\log\det\Sigma_{\mathrm{CS}}(\hat\rho_s)=d\,\mathcal{L}(\boldsymbol{\ell})$,
which is exactly $d$ times Eq~\eqref{eq:loglik} after substituting
$\hat\rho_s=(C_s-n_s)/(n_s^2-n_s)$ into the eigenvalue expression. The project's
objective is therefore the maximum log-likelihood of the model up to an additive
constant and the factor $d$, and the per-module summand diverges to $+\infty$ as
$C_s\to n_s^2$ ($\rho_s\to1$, a perfectly correlated module); we cap this at a
large finite value to keep the optimizer well-behaved.

\emph{Model selection.} The $K$-module model has $K$ free coherence parameters
over $d$ samples, giving $\mathrm{BIC}=-2d\mathcal{L}+K\log d$ and
$\mathrm{AIC}=-2d\mathcal{L}+2K$ (Eq~\eqref{eq:loglikic}). Because the likelihood
carries the factor $d$ and the penalty does not, the criteria license more
modules as the number of samples grows, as required.

\emph{Rank selection for the factor block model.} The factor model of
Eq~\eqref{eq:factorblock} has $p_m = n_m q - q(q-1)/2 + 1$ free parameters per
module, so $\mathrm{BIC}=-2d\mathcal{L}^{\text{factor}}+\log(d)\sum_m p_m$ and
$\mathrm{AIC}=-2d\mathcal{L}^{\text{factor}}+2\sum_m p_m$. Because the parameter
count depends on $q$, these are comparable across ranks and can be used to select
$q$. Evaluated at the planted partition, BIC recovers the generating rank on every
process we tested ($q=1$ for one-factor and hub designs, $q=3$ for the
three-factor design), whereas AIC selects one rank too many on the three-factor
design; we therefore recommend BIC. In practice the rank is selected by running
the optimizer once per candidate rank and keeping the lowest-BIC fit, which
attained the same recovery as the best rank chosen with knowledge of the truth on
every process.

\paragraph*{Appendix B.}
\namedlabel{AppendixB}{Appendix B}
\textbf{Reproduction package.} For convenience, a compressed version of the repository is shared as a companion to the paper (although it is also available through GitHub and Zenodo). It contains the source code, one script per analysis in \texttt{paper/experiments/} (listed in that directory's README), and the generated results tables and figures.

\section*{Data and Software Availability}
All source code for the \tool{} package, along with the reproduction scripts to generate the results and figures presented in this manuscript, is publicly available under an open-source license on GitHub at \url{https://github.com/luismadrigal98/GCM-MRK}. To ensure long-term reproducibility, a persistent snapshot of the codebase at the time of submission has been archived on Zenodo (DOI: 10.5281/zenodo.21361265). The breast-cancer RNA-seq dataset (GSE183947) analyzed in this study is publicly available from the NCBI Gene Expression Omnibus (GEO). The expression compendia and curated module definitions of the Saelens et al.\ benchmark are available from Zenodo (DOI: 10.5281/zenodo.5532578); the scripts in \texttt{paper/experiments/} document how they were subset and scored.

\section*{Acknowledgments}
We thank all the class of Simulation in Python offered in the University of Kansas, where all this project of evolving solutions using evolutionary algorithms (the simulation bit) started. We also acknowledge the use of Claude Opus 5 model from Anthropic for polishing writing and language refinement. All AI-refined content was reviewed by the authors and we take full responsibility for the content presented in the manuscript. 
